\section{Conclusion}

In this paper, we proposed a novel 4-bit BFP format, HiF4, for deep 
learning. 
Behind the design of HiF4, we first argue that E1M2 (S1P2) with 3-bit 
significand offers higher precision ceiling than E2M1 with 2-bit 
significand in the 4-bit BFP format design. 
Then, we successfully identified that increasing group size 
to 64 and employing two-level shared micro-exponents with a metadata 
overhead of 0.5-bit per value offers a new trade-off among 
overall accuracy, quantization complexity, and hardware efficiency. 
Additionally, we meticulously designed a dedicated global base scale 
that delivers sufficient inter-group dynamic range capable of directly 
supporting both training and inference scenarios, 
while enhancing the utilization of available intra-group dynamic range. 

Numerical analysis shows that, compared to the state-of-the-art 
NVFP4 format, HiF4 reduces the MSE quantization error by 24\% on  
Gaussian-distributed data. 
When 64-length dot product is integrated into existing hardware 
dot-product units originally optimized for 16-bit (FP16/BF16) 
and 8-bit (INT8/Float8) precision, 
HiF4 occupies only approximately one-third of the incremental 
area of NVFP4 and reduces power consumption by about 10\%. 
Finally, we conducted extensive LLM inference experiments, 
demonstrating that HiF4 consistently achieves higher accuracy 
than NVFP4 across various LLM architectures and diverse tasks. 
Future work will present the training potential of the proposed 
HiF4 format.



